% Generated draft: requires mathematical and source review.
\documentclass[11pt]{article}
\usepackage[margin=25mm]{geometry}
\usepackage{fontspec}
\setmainfont{Latin Modern Roman}
\usepackage{amsmath,amssymb,mathtools,unicode-math}
\setmathfont{Latin Modern Math}
\usepackage{enumitem,needspace,xurl,hyperref}
\setlength{\parindent}{0pt}
\setlength{\parskip}{5pt}
\setlength{\emergencystretch}{3em}
\newcommand{\N}{\mathbb{N}}
\newcommand{\Z}{\mathbb{Z}}
\newcommand{\Q}{\mathbb{Q}}
\newcommand{\R}{\mathbb{R}}
\newcommand{\C}{\mathbb{C}}
\newcommand{\F}{\mathbb{F}}
\newcommand{\A}{\mathbb{A}}
\newcommand{\PP}{\mathbb{P}}
\newcommand{\E}{\mathbb{E}}
\newcommand{\eps}{\varepsilon}
\newcommand{\dd}{\,\mathrm{d}}
\begin{document}
\section*{UnsolvedMath ID: 701470}
\textbf{Source problem code:} catalogue entry
\subsection{Definitions and mathematical statement}
Let \(F\) be a free group of rank \(2\). For \(x,y\in F\), write \([x,y]=xyx^{-1}y^{-1}\), and let \([F,F]\) denote the subgroup generated by these commutators. For \(g\in[F,F]\), its commutator length \(\operatorname{cl}_F(g)\) is the least nonnegative integer \(k\) such that \(g\) is a product of \(k\) commutators in \(F\), with the empty product equal to \(1\). Define its stable commutator length by
\[
\operatorname{scl}_F(g)=\lim_{n\to\infty}\frac{\operatorname{cl}_F(g^n)}{n},
\]
where \(n\) ranges over positive integers; this limit exists by subadditivity. Put
\[
S=\{\operatorname{scl}_F(g):g\in[F,F]\}\subseteq\R.
\]
The supplied source states as known background that, on this domain,
\[
\operatorname{scl}_F(g)=0\iff g=1,
\qquad
\operatorname{scl}_F(g)\ge\tfrac12\quad\text{for every }g\in[F,F]\setminus\{1\}.
\]
D. Calegari asks the following questions.
\begin{enumerate}
\item[(a)] Is \(\tfrac12\) an isolated value of \(S\)? Explicitly, does there exist \(\varepsilon>0\) such that
\[
S\cap\bigl(\tfrac12-\varepsilon,\tfrac12+\varepsilon\bigr)=\{\tfrac12\}?
\]
\item[(b)] Does there exist an interval \(J\subseteq\R\) such that \(S\cap J\) is dense in \(J\)? Here density means
\[
\forall x\in J\ \forall\varepsilon>0\ \exists g\in[F,F]:\quad
\operatorname{scl}_F(g)\in J\ \text{ and }\ 
\bigl|\operatorname{scl}_F(g)-x\bigr|<\varepsilon.
\]
The source does not specify whether \(J\) must be nonempty or have positive length. A nondegenerate-interval formulation would additionally require that \(J\) contain two distinct points. This restriction is not imposed silently here; the intended class of intervals requires clarification.
\item[(c)] Does there exist \(T\in\R\) such that
\[
\forall q\in\Q\quad
q\ge T\ \Longrightarrow\ \exists g\in[F,F]\text{ with }\operatorname{scl}_F(g)=q?
\]
\end{enumerate}
\subsection{Short English statement}
For a free group of rank \(2\), consider the stable commutator lengths of elements of its commutator subgroup. The questions ask whether \(1/2\) is isolated among these values, whether the values are dense in an interval, and whether they include every sufficiently large rational number. The intended nondegeneracy condition on the interval in part (b) is unspecified.
\subsection{Sources}
\begin{itemize}
\item Local UnsolvedMath registry, supplied complete record with permanent ID \(701470\), exact source problem code \texttt{catalogue entry}, title \emph{Kourovka Notebook Problem 18.41}; upstream revision \texttt{d371ee917ac9647fbb06cd6fc529deec6f3dda82}.
\item Supplied attribution: D. Calegari, proposed in \(2014\); original problem number \(18.41\) in \emph{The Kourovka Notebook (Unsolved Problems in Group Theory)}, 18th issue (2014). The record also identifies the 21st-issue edition (2026), \texttt{arXiv:1401.0300v48}, and names E. I. Khukhro and V. D. Mazurov as editors of the collection.
\item The local record labels the problem \emph{open} and was updated on \(2026\text{-}10\text{-}16\). This is supplied registry information, not an independently verified assessment of the literature. The cited source and edition have not been independently consulted for this draft.
\end{itemize}
\end{document}
